(1) 将随机方程写成\(dv=-av\,dt+b\,dW_t\)，\(a=\gamma/m\)、\(b=\sqrt{2\gamma k_BT}/m\)，解为\(v(t)=v(0)e^{-at}+b\int_0^te^{-a(t-s)}dW_s\)。因噪声均值零，\(\langle v(t)\rangle=\langle v(0)\rangle e^{-t/\tau}\)，\(\tau=m/\gamma\)。对\(v^2\)应用Itô公式\(d(v^2)=2v\,dv+b^2dt\)，平衡方差满足\(0=-2a\langle v^2\rangle+b^2\)，故\(\langle v^2\rangle=k_BT/m\)。
(2) 未来噪声增量独立于初速度，故上式乘\(v(0)\)取平均，\(\langle v(t)v(0)\rangle=e^{-t/\tau}\langle v(0)^2\rangle=(k_BT/m)e^{-t/\tau}\)，对负时差依平稳性为\((k_BT/m)e^{-|t|/\tau}\)。这是涨落耗散定理的时间域表现：阻尼越大，速度记忆时间越短。
(3) 位移\(\Delta x(t)=\int_0^t v(s)ds\)，所以\(\langle\Delta x^2\rangle=\int_0^t\int_0^t\langle v(s)v(s')\rangle dsds'=2(k_BT/m)\int_0^t(t-u)e^{-u/\tau}du\)。积分给\(\langle\Delta x^2\rangle=2(k_BT/\gamma)[t-\tau(1-e^{-t/\tau})]\)，此式在全时间范围从惯性连续过渡到扩散。
(4) 当\(t\ll\tau\)，括号展开为\(t^2/(2\tau)+O(t^3)\)，因此\(\langle\Delta x^2\rangle\simeq(k_BT/m)t^2\)，为弹道段；当\(t\gg\tau\)，\(\langle\Delta x^2\rangle\simeq2Dt\)，\(D=k_BT/\gamma\)。这就是Einstein关系，若噪声强度不是\(\sqrt{2\gamma k_BT}\)将无法同时得到等分配和该扩散系数。
